% FFT-based resize pipeline (resize_mode: fft).
% Adds tiny spatial- / frequency-domain thumbnails so the reader can see
% what each step does to the data.
\begin{tikzpicture}[
  node distance = 4mm and 7mm,
  every node/.style = {font=\scriptsize, inner sep=1.5pt}]

  % ---------- thumbnail helper styles --------------------------------
  % spatial-domain thumbnail: a few diagonal "wave" stripes
  % frequency-domain thumbnail: a centred bright DC bin with a faint halo

  % ---------- input spatial thumbnail --------------------------------
  \begin{scope}[local bounding box=xthumb]
    \fill[blue!12] (0,0) rectangle (0.7,0.7);
    \foreach \k in {0.05, 0.20, 0.35, 0.50, 0.65} {
      \draw[blue!60!black, thin] (0,\k) -- (0.7-\k, 0.7);
      \draw[blue!60!black, thin] (\k, 0) -- (0.7, 0.7-\k);
    }
    \draw[draw=blue!70!black, thick] (0,0) rectangle (0.7,0.7);
  \end{scope}
  \node[meta, below=1mm of xthumb] {input};
  \node[meta, above=0mm of xthumb] {$C{\times}H_n{\times}W_n$};
  % ---------- F2: full spectrum thumbnail ----------------------------
  \node[op, right=of xthumb] (fft) {$\mathcal{F}_{2}$};
  \begin{scope}[local bounding box=spec1, shift={($(fft.east) + (0.65, -0.35)$)}]
    \fill[orange!12] (0,0) rectangle (0.7,0.7);
    % bright DC bin + halo
    \foreach \r/\op in {0.30/35, 0.20/55, 0.12/75, 0.06/95} {
      \fill[orange!\op] (0.35,0.35) circle (\r);
    }
    \draw[orange!80!black, thick, dashed] (0.35,0) -- (0.35,0.7);
    \draw[orange!80!black, thick, dashed] (0,0.35) -- (0.7,0.35);
    \draw[draw=orange!80!black, thick] (0,0) rectangle (0.7,0.7);
  \end{scope}
  \node[meta, below=1mm of spec1] {fftshift};
  \node[meta, above=0mm of spec1] {DC at centre};
  \draw[arrow, thick] (xthumb.east) -- (fft.west);
  \draw[arrow, thick] (fft.east) -- ($(spec1.west) + (0,0.0)$);

  % ---------- crop / pad ---------------------------------------------
  \node[op, right=of spec1, align=center, fill=yellow!18] (cp)
    {crop\,/\,pad\\around DC};
  \begin{scope}[local bounding box=spec2,
                shift={($(cp.east) + (0.55, -0.27)$)}]
    \fill[orange!12] (0,0) rectangle (0.55,0.55);
    \foreach \r/\op in {0.22/35, 0.14/55, 0.08/75, 0.04/95} {
      \fill[orange!\op] (0.275,0.275) circle (\r);
    }
    \draw[draw=orange!80!black, thick] (0,0) rectangle (0.55,0.55);
  \end{scope}
  \draw[arrow, thick] (spec1.east) -- (cp.west);
  \draw[arrow, thick] (cp.east) -- ($(spec2.west) + (0,0)$);
  \node[meta, below=1mm of spec2] {target $H{\times}W$};

  % ---------- ifftshift, F2^-1, scale --------------------------------
  \node[op, right=of spec2] (ift) {$\mathcal{F}_{2}^{-1}$};
  \node[op, right=of ift,  align=center] (sc)  {$\times\,\rho$};
  \node[meta, below=1mm of sc] {$\rho{=}HW/H_nW_n$};

  % ---------- output spatial thumbnail -------------------------------
  \begin{scope}[local bounding box=ythumb,
                shift={($(sc.east) + (0.65, -0.35)$)}]
    \fill[blue!12] (0,0) rectangle (0.7,0.7);
    \foreach \k in {0.10, 0.30, 0.50} {
      \draw[blue!60!black, thin] (0,\k) -- (0.7-\k, 0.7);
      \draw[blue!60!black, thin] (\k, 0) -- (0.7, 0.7-\k);
    }
    \draw[draw=blue!70!black, thick] (0,0) rectangle (0.7,0.7);
  \end{scope}
  \node[meta, below=1mm of ythumb] {output};
  \node[meta, above=0mm of ythumb] {$C{\times}H{\times}W$};
  \draw[arrow, thick] (spec2.east) -- (ift.west);
  \draw[arrow, thick] (ift.east) -- (sc.west);
  \draw[arrow, thick] (sc.east) -- ($(ythumb.west) + (0,0)$);

  % bottom note
  \node[meta, below=8mm of cp, align=center, gray!70!black]
    {centre crop preserves the lowest spatial frequencies;\\
     periodic boundary conditions are exact (no interpolation, no window)};
\end{tikzpicture}
